% Tabel pendamping Gambar IV.2. Seluruh nama komponen diverifikasi langsung terhadap
% berkas sumber pada 20 Agustus 2026, bukan disalin dari catatan terdahulu.
% Lebar: 1,0 + 2,9 + 3,5 + 5,1 = 12,5 cm; tabcolsep 6pt menambah 6 x 0,21 = 1,26 cm;
% total 13,76 cm terhadap textwidth 14,00 cm.
\begin{longtable}{@{}p{1.0cm}p{2.9cm}p{3.5cm}p{5.1cm}@{}}
\caption{Pemetaan tiap tahap pada Gambar~\ref{fig:alur-rinci} ke komponen yang benar-benar dijalankan di dalam sistem. Kode pada kolom pertama sama dengan kode kotak pada gambar tersebut. Tanda $\star$ menandai tahap yang menjadi kontribusi metodologis.}
\label{tab:komponen-alur} \\
\toprule
\textbf{Kode} & \textbf{Tahap} & \textbf{Komponen dalam sistem} & \textbf{Peran} \\
\midrule
\endfirsthead

\multicolumn{4}{@{}l}{\small\emph{Tabel \thetable{} (lanjutan)}} \\
\toprule
\textbf{Kode} & \textbf{Tahap} & \textbf{Komponen dalam sistem} & \textbf{Peran} \\
\midrule
\endhead

\bottomrule
\endlastfoot

\multicolumn{4}{@{}l}{\textit{Tahap luring, dijalankan sekali di luar aplikasi}} \\
\midrule

A & Rekaman pemantauan glukosa & \texttt{data/raw/OhioT1DM/} (berkas XML) & Rekaman mentah dua belas pasien, memuat kejadian glukosa, bolus, karbohidrat, serta latihan pada stempel waktunya masing-masing. \\[2pt]

B & Penyelarasan waktu kejadian & \texttt{parse\_\allowbreak ohio\_\allowbreak xml}, \texttt{\_\allowbreak align\_\allowbreak sum} pada \texttt{ohio\_\allowbreak parser.py} & Menyelaraskan kejadian yang stempel waktunya tidak seragam ke satu kisi lima menit. \\[2pt]

C & Deret waktu terpadu & \texttt{ohio\_\allowbreak t1dm\_\allowbreak merged.csv} & Satu tabel berkolom seragam per pasien, menjadi masukan tunggal bagi seluruh tahap luring berikutnya. \\[2pt]

D & Rekayasa fitur fisiologis & \texttt{DataPreprocessor} pada \texttt{preprocessor.py} & Menurunkan tujuh fitur model, mencakup peluruhan insulin aktif dan karbohidrat aktif serta penyandian waktu dalam hari. \\[2pt]

E & Pelatihan model prakiraan & \texttt{GBMGlucoseModel} pada \texttt{gbm\_\allowbreak model.py} & Melatih model regresi untuk horizon $+30$ dan $+60$ menit, lalu menyimpannya sebagai bundel inferensi. \\[2pt]

E2 & Kalibrasi interval ketidakpastian & \texttt{conformal\_\allowbreak calibration.py}, \texttt{src/conformal.py} & Menghitung faktor konformal terpisah atas lipatan kalibrasi, sehingga cakupan 95\% menjadi terukur dan bukan diandaikan. \\[2pt]

E3 & Pelatihan pengklasifikasi kondisi~$\star$ & \texttt{train\_\allowbreak condition\_\allowbreak classifier.py} & Melatih pengklasifikasi tiga kelas sadar-biaya, sehingga kelas hipoglikemia yang jarang tidak tenggelam oleh kelas mayoritas. \\[2pt]

F & Korpus pedoman klinis & \texttt{data/knowledge\_\allowbreak base/} (berkas PDF) & Pedoman PERKENI beserta ADA sebagai satu-satunya sumber isi rekomendasi. \\[2pt]

G & Pemecahan pada batas kalimat & \texttt{MedicalKnowledgeBase} pada \texttt{knowledge\_\allowbreak base.py} & Memecah dokumen menjadi potongan yang berakhir pada batas kalimat, sekaligus melekatkan metadata identitas dokumen beserta nomor halamannya. \\[2pt]

H & Basis pengetahuan terindeks & Koleksi \texttt{diabetes\_\allowbreak kb} pada ChromaDB & Menyimpan potongan beserta metadatanya, menjadi sumber tunggal bagi ketiga cara penelusuran pada Tabel~\ref{tab:cara-penelusuran}. \\[2pt]

\midrule
\multicolumn{4}{@{}l}{\textit{Tahap daring, dijalankan pada tiap konsultasi}} \\
\midrule

I & Masukan logbook oleh dokter & \texttt{save\_\allowbreak entry} pada \texttt{1\_\allowbreak Input\_\allowbreak Logbook.py} & Merekam catatan klinis satu titik waktu ke berkas logbook. \\[2pt]

J & Pembentukan jendela dan fitur & \texttt{build\_\allowbreak window} pada \texttt{streamlit\_\allowbreak app.py}; \texttt{gabung\_\allowbreak dengan\_\allowbreak dataset}, \texttt{periksa\_\allowbreak kelayakan} pada \texttt{logbook.py} & Menyusun jendela dua belas langkah, menggabungkannya dengan deret dataset bila tersedia, lalu menolak memprediksi bila syarat kerapatan catatannya tidak terpenuhi. \\[2pt]

K & Prakiraan glukosa & \texttt{predict\_\allowbreak next} pada \texttt{streamlit\_\allowbreak app.py} & Menghasilkan nilai terprediksi $+30$ dan $+60$ menit memakai bundel hasil tahap~E. \\[2pt]

L & Interval ketidakpastian terkalibrasi & \texttt{prediction\_\allowbreak interval} pada \texttt{src/conformal.py} & Mengubah sebaran kuantil menjadi satu interval berlebar terkalibrasi yang ditampilkan kepada dokter. \\[2pt]

M & Penentuan kondisi terprediksi~$\star$ & \texttt{predict\_\allowbreak condition} pada \texttt{streamlit\_\allowbreak app.py} & Menetapkan kelas kondisi dari fitur jendela memakai pengklasifikasi hasil tahap~E3, bukan dari ambang atas nilai regresi. \\[2pt]

N & Deteksi divergensi & \texttt{evaluate\_\allowbreak divergence}, \texttt{DivergenceAlert} pada \texttt{alerts.py} & Menyalakan peringatan ketika kondisi terkini berbeda dari kondisi terprediksi, yakni tepat pada keadaan saat pengondisian prediksi berpengaruh. \\[2pt]

O & Perakitan kondisi klinis terstruktur~$\star$ & \texttt{PatientState} pada \texttt{patient\_\allowbreak state.py} & Merangkum nilai terprediksi, kelas kondisi, arah tren, serta konteks fisiologis menjadi satu objek yang menjadi masukan tunggal bagi pembentuk kueri. \\[2pt]

P & Transformasi kueri terkondisi-prediksi~$\star$ & \texttt{PredictionConditionedQueryBuilder} pada \texttt{conditioned\_\allowbreak query.py} & Menyusun kueri dari kondisi yang akan datang alih-alih dari kondisi terkini. Inilah titik kebaruan penelitian ini. \\[2pt]

Q & Penelusuran leksikal & \texttt{SimpleKeywordRetriever} pada \texttt{retriever.py} & Memilih lima potongan teratas dari basis tahap~H memakai pembobotan Okapi BM25. \\[2pt]

R & Pembangkitan rekomendasi & \texttt{RAGGenerator}, \texttt{DiabetesAdvisorChain}, \texttt{RAGPipeline} & Menyusun rekomendasi berbahasa Indonesia yang dibumikan pada kelima potongan tersebut, tanpa pernah menerima nomor halaman. \\[2pt]

S & Penjaminan penyangkalan & \texttt{prompts.py} beserta \texttt{RAGPipeline} & Memastikan tiap keluaran memuat pernyataan bahwa keputusan klinis akhir berada pada dokter. \\[2pt]

T & Resolusi sitasi dari metadata~$\star$ & \texttt{build\_\allowbreak source\_\allowbreak list}, \texttt{format\_\allowbreak page\_\allowbreak label} pada \texttt{citations.py} & Menyusun daftar sumber beserta nomor halamannya dari metadata potongan, sehingga nomor halaman tidak mungkin dikarang model bahasa. \\[2pt]

U & Layar konsultasi dokter & \texttt{streamlit\_\allowbreak app.py} beserta \texttt{app/ui.py} & Menyajikan prakiraan, interval, peringatan, rekomendasi, serta potongan sumbernya pada satu layar. \\[2pt]

V & Pencatatan keputusan dokter & \texttt{ClinicalDecisionLog} pada \texttt{decision\_\allowbreak log.py} & Menyimpan keputusan beserta potongan yang persis dilihat dokter, sehingga keputusan tersebut dapat diaudit di kemudian hari. \\

\end{longtable}

{\footnotesize Tanda $\star$ muncul pada \textbf{lima} baris untuk \textbf{empat} tahap kontribusi, sebab pengklasifikasi kondisi tercatat dua kali: tahap~E3 adalah pelatihannya di sisi luring, sedangkan tahap~M adalah pemakaiannya di sisi daring. Keduanya satu komponen yang sama.\par}
